\section{潜在结果把处理效应定义在替代世界之间}
\label{sec:13-Machine-Learning/12-Causal-Inference/04}

一种新药是否降低血压？你找来 200 个受试者，100 人服药，100 人服安慰剂。四周后，服药组平均收缩压降了 8 mmHg，对照组降了 2 mmHg。差值 6 mmHg——能归因于药吗？

你心里清楚：如果当初选择服药的那 100 人和选择安慰剂的那 100 人，在年龄、基础血压、运动习惯上本来就不同，那么 6 mmHg 也许只是这些先存差异的投影，和药效毫无关系。

更深的麻烦在后面。即使两组在已测变量上完全匹配，你也永远不会知道：某个服了药的受试者，如果当初没有服药，他的血压会是多少。你只看到了他所在的那条世界线上发生的事。替代世界线——同一个体在另一种处理下的结果——是永远缺失的数据。

本节要讲的就是：怎样在永远缺失半边数据的条件下，仍然定义和识别处理效应。答案的核心是潜在结果记号：先把每个个体在每种可能处理下的结果定义出来，承认它们不能同时观察，然后用设计或识别假设把不同个体的可观察部分拼接成总体比较。

\subsection{潜在结果：先定义个体在两个世界下的值}

对每个单位 \(i\)，定义两个量：

\[
Y_i(1) = \text{若 } i \text{ 接受处理时的结果},\qquad
Y_i(0) = \text{若 } i \text{ 不接受处理时的结果}.
\]

个体的处理效应就是二者之差：

\[
Y_i(1) - Y_i(0).
\]

这个定义看起来再直白不过。但它立刻制造了一个无法回避的困境：同一个单位在同一时刻只能经历一种处理。\(Y_i(1)\) 和 \(Y_i(0)\) 中，你最多只能观察到一个。另一个永远是反事实的。

这不是样本量大了就会消失的缺失值问题。你不可能让同一个人在同一个四周里既服药又不服药然后比较——这不是数据采集的成本问题，是物理上不可能。Holland (1986) 称之为因果推断的根本问题。

潜在结果框架的价值在于：它先把``处理效应''定义得干干净净，然后坦白承认其中一半不可观察。接下来的全部论证——随机化、可忽略性、工具变量——都不是为了绕开这个根本问题，而是为了在它面前找到有意义的总体量仍然可以识别。

\subsection{一致性：把观察结果拴在潜在结果上}

你还需要一座桥，把现实世界中实际观察到的结果 \(Y_i\) 和那两个潜在世界中的 \(Y_i(1), Y_i(0)\) 连接起来。令 \(X_i \in \{0, 1\}\) 表示实际接受的处理，一致性公理要求：

\[
Y_i = X_i Y_i(1) + (1 - X_i) Y_i(0).
\]

这条式子说了两件事。第一，某个个体实际接受的是处理 1 时，你观察到的就是他潜在结果 \(Y_i(1)\)。第二——也是更容易被忽略的——``处理 1''必须有足够明确且稳定的版本：同一标签下的不同剂量、不同执行质量或不同辅助照护，不能混为一谈。如果把 10 mg 和 50 mg、以及附带营养咨询和不附带的干预全部标成``处理 1''，一致性就只是记号游戏，不描述任何真实的潜在结果。

一致性公理的另一面是通常要求的稳定单元处理值假设（SUTVA）。SUTVA 由两部分组成：一部分是一致性本身；另一部分是 no-interference——一个单位的潜在结果不受其他单位处理状态的影响。

No-interference 在实际中容易违反的场景需要被正视。传染病传播——你的感染概率依赖周围人的接种覆盖率，而非仅仅你自己的接种状态；社交网络效应——你的行为受朋友是否接受干预影响；班级教学——你的成绩依赖整个班级的分组方案，而非只取决于自己被分到哪一组；市场竞争——一个区域的定价效果受相邻区域定价策略影响。在这些场景中，潜在结果应写成全体处理向量的函数 \(Y_i(X_1, \ldots, X_n)\)，而不是只依赖自己的 \(X_i\)。这种一般化大幅增加了维度，识别通常需要额外的网络结构或空间结构假设来约束干扰的传播范围。

\subsection{平均效应可以识别，个体缺失世界仍不会回来}

由于个体的 \(Y_i(1) - Y_i(0)\) 无法直接计算，因果推断把目标转向总体平均量。最基础的目标是总体平均处理效应：

\[
\tau = \E[Y(1) - Y(0)].
\]

依应用场景不同，还有几个变体。已处理人群平均效应衡量实际接受了处理的那批人受益多少：

\[
\tau_{\mathrm{ATT}} = \E[Y(1) - Y(0) \mid X = 1].
\]

条件平均处理效应在给定协变量 \(C = c\) 的子群体中估计效应：

\[
\tau(c) = \E[Y(1) - Y(0) \mid C = c].
\]

当处理效应因人而异时，ATE、ATT 和 CATE 可以是三个完全不同的数。举个具体的数例：假设人群由两类人组成，各占一半。第一类人服药降 10 mmHg，第二类人服药降 0 mmHg。如果选择服药的人中第一类占 80\%，那么 ATT \(= 0.8 \times 10 + 0.2 \times 0 = 8\) mmHg，而 ATE \(= 0.5 \times 10 + 0.5 \times 0 = 5\) mmHg。ATT 大于 ATE——因为从治疗中获益更多的人更倾向于选择治疗。这个差异不是偏倚，它是两个不同的因果量在说实话。

这些总体平均量能在适当设计或假设下由不同单位的可观察部分拼出来。但你要清楚：即使 \(\tau\) 被完美识别，它也不会告诉你任何一个具体个体的 \(Y_i(1) - Y_i(0)\)。模型输出的个体效应预测——无论输出到小数点后几位——仍含有不可验证的跨世界外推，不是观察事实。

\subsection{随机分配切断处理与潜在结果的联系}

在理想随机试验中，谁进处理组由抛硬币决定，和这个人本身的所有特征——包括他的潜在结果——无关：

\[
X \perp \bigl(Y(1), Y(0)\bigr).
\]

有了这个独立性和一致性，最简单的组间均值差就识别了 ATE：

\[
\begin{aligned}
\E[Y \mid X = 1] - \E[Y \mid X = 0]
&= \E[Y(1) \mid X = 1] - \E[Y(0) \mid X = 0] \\
&= \E[Y(1)] - \E[Y(0)].
\end{aligned}
\]

第一步用了一致性，第二步用了独立性。两个等号的跨度看起来不大，但它们跨越了从可观察量到潜在结果、从条件期望到边缘期望的两道门槛。

这里的识别力量来自已知的随机化机制，不是来自两组基线变量在有限样本中``恰好看起来接近''。有限样本不平衡不会破坏识别的无偏性，但会影响精度；更严重的是执行层面的偏差——受试者退出、不依从方案、失访——这些会改变实际被估计的对象。

按分配组比较得到的是 intention-to-treat effect（ITT），保持了随机化的保护。ITT 回答的问题是``邀请人们接受治疗的因果效应''而不是``实际接受治疗的因果效应''。只比较实际接受了处理的人则重新引入了选择——那些选择依从的人可能和选择不依从的人在潜在结果上本就不同。如果目标是 treatment-on-the-treated，需要关于依从类型、排除限制或工具变量的额外结构。工具变量框架正是为这一情形设计的：随机分配 \(Z\) 充当工具——它影响实际治疗 \(X\) 但不直接影响结果（排除限制），而它影响结果的唯一途径就是通过 \(X\)。对某些依从子群体（compliers），实际接受的处理的因果效应是可以识别的。

\subsection{观察研究用条件可忽略性替代随机化}

没有随机化时，处理分配可能和潜在结果相关——病情重的人更可能接受积极治疗，他们即使不接受治疗结果本来就更差。此时 \(\E[Y \mid X = 1] - \E[Y \mid X = 0]\) 混合了处理效应和选择偏差，无法直接解释为因果效应。

如果有足够丰富的一组处理前协变量 \(C\)，使得在 \(C\) 的每一层内部，处理分配近似于随机：

\[
\bigl(Y(1), Y(0)\bigr) \perp X \mid C,
\]

则称在给定 \(C\) 后处理分配可忽略。再配上重叠条件

\[
0 < P(X = 1 \mid C = c) < 1,
\]

即对每一类 \(c\)，两种处理都实际有出现，就可以通过分层先在各层内部估计效应再整合：

\[
\E[Y(x)] = \E_C\!\left[\E[Y \mid X = x, C]\right].
\]

这个识别公式与后门调整公式是同一个结论的两种表达——一个用潜在结果语言，一个用图模型语言。它们都在回答：在什么条件下，条件关联可以被提升为因果效应。

条件可忽略性不是拟合优度能检验出来的。它声明 \(C\) 已经囊括了处理与潜在结果的全部共同原因——这是一个关于不可观察世界的断言，数据中的任何模式都无法直接确认它。你可以检查的只是：加入 \(C\) 后处理效应的估计是否对 \(C\) 的进一步细化保持稳定；但这也不能证明没有未测量的混杂，只能说明在已测量混杂的尺度上估计是稳健的。

重叠条件同样关键。如果某类重症患者必然接受处理，则该层内没有未处理者提供反事实比较——倾向得分在该层为 1。此时识别在无数据区域只能靠模型外推。外推的质量取决于模型假设的形式（线性、多项式、样条），不取决于数据对该区域的提示——因为根本没有数据提示。倾向得分接近 0 或 1 时，逆概率加权的方差急剧膨胀，实际操作中通常需要截断（trimming）或使用重叠加权。

对于无法检验可忽略性的困境，敏感性分析提供了一种正视之道：假设存在一个未测量的混杂因子 \(U\)，设它与处理和结果的关联强度为某参数 \(\gamma\)，然后看 \(\gamma\) 需要大到什么程度才能把估计的处理效应完全解释掉。如果只需要很弱的未测混杂就能把效应归零，结论就是脆弱的；如果需要强到不合理的混杂，结论就更可信。这不是在检验假设是否成立，而是在量化结论对假设违反的耐受程度。

\subsection{设计先于估计}

随机化、自然实验、断点回归、工具变量——这些设计的共同目标是从分配机制中获得可信的比较，而不是先选一个回归模型再解释系数。

观察研究中的匹配、分层和逆概率加权，做的也是同一件事的近似版本：在已测协变量上重建一个可比的伪总体，使得处理组和对照组在 \(C\) 的分布上尽可能接近。倾向得分在此处的角色是降维工具：它不是因果量，而是一个平衡性得分——在倾向得分的每一层内，处理组和对照组在 \(C\) 上的分布是平衡的。这样就把多维协变量的匹配问题压缩成了一维倾向得分的匹配或加权问题。

一个实际的研究流程也许比任何公式都重要：在研究开始时就固定处理的版本（什么剂量、多长时间）、结果的时间窗、目标人群和 estimand——是 ATE、ATT 还是某个条件效应。如果这些定义随着看到结果后再调整，后面计算的标准误再精确，也只是为移动后的新问题提供精确数字，而不是为最初的问题提供可靠答案。Gelman 和 Loken 用``garden of forking paths''来形容这种现象：同一个人面对同一个数据集，可以通过无数种合法的分析选择到达截然不同的结论，而 p 值没有为这种自由度留出折扣。

潜在结果框架不负责替你判断因果——它负责把因果问题写成可以被设计、被识别假设和被数据联合回答的形式。它明确告诉你：哪个量来自观察、哪个量来自假设、哪个量永远不可知。分得清这三层，才是因果推断真正的起点。
